Los requisitos funcionales del \dfn{backend} se agrupan en 16 subsistemas
que totalizan 131 sub-requisitos, elicitados a partir de las reuniones con
el tutor, el estudio comparativo de plataformas del
\cref{CAP:STATE-OF-THE-ART} y los requisitos institucionales de seguridad
y privacidad. La \cref{table:rf-summary} resume los grupos y su conteo;
a continuación se describe la responsabilidad de cada uno. La
especificación detallada de cada sub-requisito, incluyendo entradas,
proceso y salida, se recoge en el \cref{AP:REQUIREMENTS}.

% --- Tabla resumen (se mantiene exactamente igual) ---
\begin{table}[Resumen de requisitos por grupo]%
            {table:rf-summary}%
            {Resumen de requisitos funcionales por grupo, con conteo de
             sub-requisitos.}
  \small
  \renewcommand{\arraystretch}{1.2}
  \begin{tabular}{c l r}
    \hline
    \textbf{N.\textsuperscript{o}} & \textbf{Subsistema} & \textbf{Sub-RF} \\
    \hline
    1  & Autenticación y sesiones                       & 5   \\
    2  & Gestión de usuarios y organizaciones           & 13  \\
    3  & Cursos académicos                              & 5   \\
    4  & Asignaturas y grupos                           & 10  \\
    5  & Permisos y autorización                        & 7   \\
    6  & Exámenes y modelos                             & 16  \\
    7  & Instancias de examen y páginas                 & 15  \\
    8  & Reglas de corrección                           & 5   \\
    9  & Ingesta de documentos y reconocimiento         & 11  \\
    10 & Anotaciones                                    & 7   \\
    11 & Calificaciones                                 & 6   \\
    12 & Revisión de exámenes                           & 8   \\
    13 & Notificaciones                                 & 7   \\
    14 & Exportación y búsqueda                         & 5   \\
    15 & Configuración y mantenimiento                  & 5   \\
    16 & Seguridad y auditoría                          & 6   \\
    \hline
       & \textbf{Total sub-requisitos funcionales}      & \textbf{131} \\
    \hline
  \end{tabular}
\end{table}

\begin{functional}

  \item \textbf{Autenticación y sesiones.}
  \label{rf:auth}
  El sistema emplea una arquitectura de autenticación basada en
  \dfnpl{plugin} intercambiables que desacoplan el mecanismo concreto
  del resto de la aplicación. En la versión~1.0 se implementa un
  \dfn{plugin} \acs{jwt} que emite pares de \textit{tokens} (acceso de
  corta duración y refresco rotativo) a partir de credenciales de correo
  y contraseña, y permite la renovación y el cierre de sesión mediante
  lista negra en Redis con \acs{ttl}. La arquitectura queda preparada
  para incorporar \acs{sso} institucional (\acs{saml}~2.0 y \acs{oidc})
  sin modificar el \dfn{frontend} ni la lógica de negocio.

  \item \textbf{Gestión de usuarios y organizaciones.}
  \label{rf:users}
  El sistema se organiza en torno a organizaciones como contenedor
  \dfn{multi-tenant}. Un superadministrador global crea las
  organizaciones con su gestor inicial; los gestores administran
  usuarios dentro de su organización (alta individual y masiva desde
  \acs{csv} o \acs{json}, modificación, desactivación lógica, gestión
  de permisos de gestor y restablecimiento de contraseñas). Cada
  usuario puede consultar su perfil, modificar su correo y gestionar
  sus preferencias de notificación. El \acs{dni} se cifra con
  \acs{aes}-256-\acs{gcm} y las contraseñas se generan automáticamente
  con requisitos de complejidad, se almacenan como \textit{hash}
  Argon2 y se envían por correo asíncrono.

  \item \textbf{Cursos académicos.}
  \label{rf:courses}
  El sistema se articula en torno a cursos académicos, con exactamente
  un curso activo por organización. La creación de un nuevo curso
  dispara una transición transaccional que archiva el curso anterior,
  desactiva a los usuarios no gestores, archiva membresías y
  notificaciones, y archiva las instancias de examen. Los gestores
  pueden navegar por cursos archivados en modo lectura y modificar los
  datos del curso activo.

  \item \textbf{Asignaturas y grupos.}
  \label{rf:subjects}
  Las asignaturas pertenecen a un curso y contienen grupos, un
  coordinador y una lista de miembros con roles (profesor o alumno). El
  coordinador o el gestor pueden crear, modificar y eliminar asignaturas
  y grupos, y gestionar las membresías. Se soporta la importación y
  exportación masiva desde \acs{csv} o \acs{json}, y cada usuario puede
  consultar las asignaturas del curso activo en las que participa con su
  rol y grupo.

  \item \textbf{Permisos y autorización.}
  \label{rf:perms}
  El sistema de autorización opera en tres capas: el \textit{flag}
  \texttt{is\_staff} otorga permisos plenos de gestor sobre todas las
  asignaturas de la organización; (\texttt{SubjectMembership}),
  la membresía de asignatura, define un rol (\texttt{COORDINATOR},
  \texttt{TEACHER} o \texttt{STUDENT}) con un conjunto de trece permisos
  contextuales por defecto, recogidos en la
  \cref{table:perms-default}; y las reglas de asignación
  (\texttt{AssignmentRule}) restringen el acceso a instancias o
  problemas concretos. Los permisos por defecto son configurables por
  organización y el coordinador puede modificar los de un profesor
  individual mediante anulaciones, con la restricción de no poder
  delegar un permiso que él mismo no posea.

  \begin{table}[Permisos por defecto por rol]{table:perms-default}{%
      Permisos contextuales y sus valores por defecto para cada rol en una
      asignatura. Un valor \texttt{Sí} indica que el rol posee el permiso
      por defecto; el coordinador puede modificar los del profesor mediante
      anulaciones. Los permisos del estudiante son fijos y no admiten
      cambios por organización ni por miembro.}
    \small
    \renewcommand{\arraystretch}{1.2}
    \begin{tabular}{l c c c}
      \hline
      \textbf{Permiso} & \textbf{COORDINATOR} & \textbf{TEACHER} & \textbf{STUDENT} \\
      \hline
      \texttt{can\_create\_exams}       & Sí & No  & No \\
      \texttt{can\_create\_rubric}      & Sí & Sí  & No \\
      \texttt{can\_assign\_correctors}  & Sí & No  & No \\
      \texttt{can\_resolve\_issues}     & Sí & No  & No \\
      \texttt{can\_publish\_grades}     & Sí & No  & No \\
      \texttt{can\_manage\_reviews}     & Sí & No  & No \\
      \texttt{can\_view\_all\_instances}& Sí & Sí  & No \\
      \texttt{can\_export\_grades}      & Sí & No  & No \\
      \texttt{can\_manage\_members}     & Sí & No  & No \\
      \texttt{can\_delegate\_perms}     & Sí & No  & No \\
      \texttt{can\_manage\_groups}      & Sí & No  & No \\
      \texttt{can\_grade}               & Sí & Sí  & No \\
      \texttt{can\_annotate}            & Sí & Sí  & No \\
      \hline
    \end{tabular}
  \end{table}

  \item \textbf{Exámenes y modelos.}
  \label{rf:exams}
  Un examen pertenece a una asignatura y se compone de uno o más
  modelos estructurales, cada uno con un \acs{pdf} en blanco de
  referencia y perfiles de página que agrupan zonas de reconocimiento
  tipadas (\texttt{OCR\_TEXT}, \texttt{OCR\_NUMBER}, \texttt{QR},
  \texttt{CHECKBOX}). El sistema permite definir problemas con
  puntuación máxima (positiva o negativa), rúbricas de
  \textit{checkmarks} con puntuación fija, y la lista de alumnos
  convocados. Sobre el \acs{pdf} en blanco se genera un \acs{pdf}
  instrumentado con códigos \acs{qr} que codifican la organización, el
  examen, el modelo y el número de página. El examen incluye permisos
  configurables que controlan qué información ve el alumno sobre su
  instancia una vez publicada.

  \item \textbf{Instancias de examen y páginas.}
  \label{rf:instances}
  Una instancia materializa el examen para un alumno concreto. Las
  páginas escaneadas (\texttt{ExamPage}) se almacenan individualmente
  en MinIO y se componen bajo demanda en un \acs{pdf} completo. Cada
  instancia sigue una máquina de estados con diez estados y
  transiciones validadas, incluyendo retrocesos autorizados. El sistema
  detecta automáticamente incidencias y permite su resolución
  manual mediante herramientas de gestión de páginas (reordenar,
  eliminar, mover, adjuntar huérfanas). La publicación masiva transita
  todas las instancias corregidas y sin incidencias al estado
  \texttt{PUBLISHED}, notificando opcionalmente a los alumnos.
  \begin{table}[Tipos de incidencia]{table:instance-issues}{%
      Incidencias que el sistema detecta durante la ingesta y el
      reconocimiento.}
    \small
    \renewcommand{\arraystretch}{1.2}
    \begin{tabular}{l p{8.5cm}}
      \hline
      \textbf{Incidencia} & \textbf{Causa} \\
      \hline
      \texttt{QR\_READ\_ERROR}       & QR no legible en la página \\
      \texttt{QR\_MISMATCH}          & QR indica examen/modelo/página incorrectos \\
      \texttt{MISSING\_PAGE}         & Falta una o más páginas tras el tiempo de espera \\
      \texttt{EXTRA\_PAGE}           & Página con número superior al máximo del modelo \\
      \texttt{ORPHAN\_PAGE}          & Página sin QR, asignada por proximidad temporal \\
      \texttt{STUDENT\_NOT\_IDENTIFIED} & Atributos de identificación sin coincidencia \\
      \texttt{DUPLICATE\_STUDENT}    & El alumno ya tiene otra instancia en el examen \\
      \hline
    \end{tabular}
  \end{table}
  Las zonas de reconocimiento se clasifican en los tipos recogidos en la
  \cref{table:zone-types}.

  \begin{table}[Tipos de zona de reconocimiento]{table:zone-types}{%
      Tipos de zona que pueden definirse en un perfil de página.}
    \small
    \renewcommand{\arraystretch}{1.2}
    \begin{tabular}{l l p{7cm}}
      \hline
      \textbf{Tipo} & \textbf{Reconocedor} & \textbf{Descripción} \\
      \hline
      \texttt{OCR\_TEXT}    & Motor \acs{ocr} & Texto libre (nombre, apellidos) \\
      \texttt{OCR\_NUMBER}  & Motor \acs{ocr} & Texto numérico (NIA, DNI) \\
      \texttt{QR}           & Decodificador fijo & Código QR con
                               \textit{checksum} de integridad \\
      \texttt{CHECKBOX}     & Densidad de píxeles & Detección binaria
                               (marcado / no marcado) \\
      \hline
    \end{tabular}
  \end{table}
  \item \textbf{Reglas de corrección.}
  \label{rf:assignment}
  El coordinador define reglas de asignación (\texttt{AssignmentRule})
  que asignan correctores a instancias o problemas mediante filtros
  combinables de grupo, modelo y problema. Las reglas se evalúan contra
  las instancias existentes para generar las asignaciones concretas,
  también al crearse nuevas instancias. El sistema verifica la
  cobertura de corrección y ofrece a cada corrector un listado de sus
  tareas pendientes.

  \item \textbf{Ingesta de documentos y reconocimiento.}
  \label{rf:ingest}
  Un \textit{watcher} monitoriza carpetas \acs{sftp} en busca de nuevas
  páginas, las valida, las almacena en MinIO y encola tareas de
  reconocimiento en Celery. Un \textit{dispatcher} ejecuta el
  reconocedor adecuado para cada zona del perfil de página:
  decodificador \acs{qr} con validación de \textit{checksum}, motor
  \acs{ocr} configurable por organización para zonas de texto y
  numéricas, y análisis de densidad de píxeles para casillas de
  verificación. Los textos extraídos se asocian a los alumnos
  convocados mediante un \textit{matching} tolerante a errores comunes
  de \acs{ocr} (confusiones entre caracteres similares como \texttt{O}
  y \texttt{0}, o \texttt{l} y \texttt{1}), ponderando las evidencias
  de todas las páginas del examen y seleccionando la mejor coincidencia
  global. Las páginas se ensamblan en instancias, transitando
  automáticamente a \texttt{RECEIVED} al completarse. Se soporta además
  \acs{ocr} sobre anotaciones de calificación y de \textit{stylus}.

  \item \textbf{Anotaciones.}
  \label{rf:annotations}
  Los correctores crean anotaciones de texto enriquecido, trazos de
  \textit{stylus} o audio sobre las páginas de una instancia. Las
  anotaciones pueden asociarse opcionalmente a un problema concreto y
  su contenido binario se almacena en MinIO. El sistema valida la
  duración de los audios y permite al autor modificar o eliminar sus
  anotaciones.

  \item \textbf{Calificaciones.}
  \label{rf:grading}
  Cada problema de una instancia recibe una calificación numérica
  decimal, que puede ser negativa a nivel de problema (la nota total de
  la instancia tiene mínimo cero). El corrector puede asignar la nota
  manualmente o marcando criterios de una rúbrica, cuyas puntuaciones
  positivas y negativas se suman para obtener la calificación del
  problema. La concurrencia se controla mediante comparación del valor
  anterior: si otro corrector modificó la nota en el ínterin, la
  operación se rechaza con \acs{http}~409 y se devuelve el valor actual
  para que el corrector pueda sincronizarse. La nota total se recalcula
  automáticamente tras cada cambio. El alumno puede consultar su nota
  una vez publicada, con el nivel de detalle que el examen permita.

  \item \textbf{Revisión de exámenes.}
  \label{rf:review}
  El coordinador define una ventana de revisión con fechas de inicio y
  fin. A la hora de inicio, las instancias publicadas transitan
  automáticamente a \texttt{IN\_REVIEW}; durante la ventana, los
  alumnos pueden solicitar la revisión de problemas concretos. Al
  cierre, las instancias sin solicitudes pasan a \texttt{FINALIZED} y
  las que tienen solicitudes quedan en \texttt{PENDING\_REVIEW}. Los
  correctores reciben una notificación agrupada con sus solicitudes
  pendientes, las resuelven y, cuando todas las de una instancia están
  cerradas, ésta transita a \texttt{FINALIZED} y se notifica al alumno
  el resultado.

  \item \textbf{Notificaciones.}
  \label{rf:notif}
  El sistema genera notificaciones \textit{in-app} para los eventos
  relevantes (publicación de notas, apertura de revisión, asignación de
  corrección, incidencias de ingesta, etc.). Si el usuario tiene
  habilitada la preferencia, se envía además un correo electrónico
  espejo mediante tareas asíncronas. Las notificaciones pueden marcarse
  como leídas individual o masivamente, se listan con paginación y
  filtros, y se archivan automáticamente durante la transición de
  curso. Los correos de bienvenida y restablecimiento de contraseña se
  envían como correos directos, sin generar notificación
  \textit{in-app}.

  \item \textbf{Exportación y búsqueda.}
  \label{rf:export}
  El sistema permite exportar las calificaciones de un examen en
  \acs{csv} (con columnas configurables por organización) y en
  \acs{json} con desglose por problema. Los gestores disponen de una
  búsqueda jerárquica para explorar el histórico de la organización
  (cursos $\to$ asignaturas $\to$ exámenes $\to$ instancias) con
  filtros combinables y búsqueda parcial por texto. Los datos de cursos
  archivados permanecen accesibles en modo solo lectura.

  \item \textbf{Configuración y mantenimiento.}
  \label{rf:config}
  Cada organización dispone de un conjunto de parámetros configurables
  en caliente (motor \acs{ocr}, umbrales de confianza, tiempos de
  espera, límites de tamaño y duración, frecuencia de borrado
  automático, idioma, permisos por defecto de los roles). Los cambios
  invalidan la caché Redis asociada. Una tarea periódica elimina los
  exámenes archivados que excedan la antigüedad configurada, notificando
  previamente a los gestores.

  \item \textbf{Seguridad y auditoría.}
  \label{rf:security}
  Un registro de auditoría \textit{append-only}, protegido a nivel de
  base de datos mediante un disparador que impide modificaciones,
  almacena cada evento de negocio con actor, organización, dirección
  \acs{ip} y datos relevantes. Los gestores pueden consultar el log con
  filtros avanzados. Los datos sensibles se protegen con cifrado en
  reposo (\acs{aes}-256-\acs{gcm} para el \acs{dni}, con \textit{nonce}
  único por registro), \textit{hashing} robusto de contraseñas (Argon2,
  diseñado para ser resistente a ataques de fuerza bruta con hardware
  especializado), marca de agua dinámica con identidad del alumno y
  fecha sobre las páginas servidas, \acs{url} temporales con expiración
  corta y cifrado del \acs{pdf} en tránsito con una clave derivada del
  identificador de sesión.

\end{functional}
